\section{Verification of the scalar estimates}\label{cert:verification}

This appendix proves the Gaussian estimates used in the transport and scalar
arguments. All terminating decimals in this appendix denote exact rational
numbers. The finite certificates below concern entire intervals: their
validity follows from coefficient bounds and polynomial positivity, with no
sampling assumption.
In the negative tail, $k$ and its first three derivatives are small and
can be approximated by Gaussian excess moments. In the positive tail,
$k\sim a/2$, so we subtract this leading term and control derivatives of
the remaining logarithmic correction. On the compact interval, we
propagate the differential equations in their stable directions---backward
from $q(10)$ and forward from $k(-6)$---and then certify polynomial
positivity by exact arithmetic.

\begin{proof}[Proof of Lemma~\ref{sc:bounds} and Lemma~\ref{tr:scalar-basic}]
We use the definitions of $q,k,d,b,l$ from the transport section and of
$p,h_0,\mathsf D,\mathsf L,\mathsf E,K_0,J_0,A_*,\eta_*,\nu_*$
from the scalar section. The differential identities needed below are
\begin{equation}\label{cert:identities}
 q'=qd,\qquad k'=b=1-kd,\qquad
 l=k-q+ab,\qquad l'=2b-qd+al.
\end{equation}
In particular,
\begin{equation}\label{cert:derived}
 \begin{split}
 p&=b+kl-l^2/\alpha,\\
 \mathsf E&=p+b(.75-b)-l^2/\alpha,\\
 h_0&=k(b-.5)-2(b-.28)l/\alpha,\\
 \frac{d\log M}{d\log q}
 &=\frac{2(1-qd)}{d^2}
       -\frac{2bl}{d(1-b^2)},\qquad M=d^{-2}(1-b^2).
 \end{split}
\end{equation}
We verify the estimates separately on $a\le-6$, $a\ge10$, and
$[-6,10]$. The central calculation is also specified in
\texttt{verification/verify\_scalar.py}; it uses rational arithmetic only.

\subsection*{The negative tail}
Put $r=-a\ge6$, $u=r^{-2}$, and
\[
 I_j(r)=\int_0^\infty v^j e^{-rv-v^2/2}\,dv\qquad(j\ge0).
\]
Writing $\varphi$ for the standard Gaussian density, the lower Gaussian tail
is $P=\Phi(-r)=\varphi(r)I_0(r)$. Thus
\[
 q=\frac{\varphi(r)}{1-P},\qquad
 k=\frac{(1-P)(-\log(1-P))}{\varphi(r)}.
\]
Integration by parts gives
\[
 I_1=1-rI_0,\qquad I_2=I_0-rI_1,\qquad I_3=2I_1-rI_2.
\]
The inequality $0\le1-(1-P)(-\log(1-P))/P\le2P$ gives
$|k-I_0|\le2P/r$. Since $P\le\varphi(r)/r$ and
$\varphi(6)<7\cdot10^{-9}$, we have $q\le1.01\varphi(r)$.
Applying \eqref{cert:identities} successively therefore gives
\[
 |b-I_1|\le4\varphi(r)/r,\quad
 |l-I_2|\le6\varphi(r),\quad
 |l'-I_3|\le9r\varphi(r).
\]
Each Gaussian error decreases faster than its required inverse power for
$r\ge6$: indeed $r^N\varphi(r)$ is decreasing there for $N\le7$.
Checking at $r=6$ yields, with $g_0=k,g_1=b,g_2=l,g_3=l'$,
\begin{equation}\label{cert:left-errors}
 |g_j-I_j|\le .01u\frac{j!}{r^{j+1}}\qquad(0\le j\le3).
\end{equation}
The elementary bounds $1-v^2/2\le e^{-v^2/2}\le1$ also give
\begin{equation}\label{cert:integral-errors}
 1-\frac{(j+1)(j+2)}2u
 \le \frac{I_j}{j!/r^{j+1}}\le1.
\end{equation}
For $\zeta=1+.01u$, consequences we shall use are
\[
 \begin{array}{c}
 0<k\le\zeta/r,\quad
 u(1-3.01u)\le b\le u\zeta<.028,\\
 2r^{-3}(1-6.01u)\le l\le2\zeta r^{-3},\quad
 0<l'\le6\zeta r^{-4}.
 \end{array}
\]
In particular $0<b<1/2$ and $l>0$. Moreover,
\[
 p\ge u(1-3.01u)-4\zeta^2u^3/\alpha
       \ge u(1-4.54u)\ge\frac1{7+r^2}.
\]
The last two inequalities follow by inserting $u\le1/36$; after dividing
by positive factors they reduce to
$3.01+4\zeta^2u/\alpha<4.54$ and
$(1-4.54u)(1+7u)\ge1$.
The two summands of $h_0$ in \eqref{cert:derived} have opposite signs
and magnitudes at most $.501/r$ and $.43/r$. Hence
\[
 h_0^2\le .501^2u<.67^2u(1-4.54u)\le .67^2p.
\]
This establishes the required bounds on $p,h_0$ in this tail.

For clarity, the remaining substitutions can be made with the following
ratios, avoiding cancellation in $p$:
\begin{equation}\label{cert:left-ratios}
 \frac{kl}{b}\le\frac{2u\zeta^2}{1-3.01u},\qquad
 \frac{l^2}{\alpha b}\le\frac{4u^2\zeta^2}{\alpha(1-3.01u)},
 \qquad .93<\frac bp<1.05.
\end{equation}
They imply the $O$ row of the coefficient table in Lemma~\ref{sc:bounds}.
For example, $\mathsf D\le D_*(7+r^2)<.295(7+r^2)$,
$\mathsf L\le2\zeta/(\alpha r^3)<.60$,
\[
 .37<K_0=D_*\left(1+( .75-b)\frac bp-
              \frac{l^2}{\alpha b}\frac bp\right)<.56,
 \qquad J_0=D_*c\frac bp<.72.
\]
The bounds for $z=l'/l$, $h_*=1+3b+kz$, and $\chi=l/d$ follow from
\[
 0<z\le\frac{3\zeta}{r(1-6.01u)}<.61,\qquad
 1<h_*<1+3(.028)+\zeta(.61)/6<1.77,
 \qquad \chi\le2\zeta u^2<.070,
\]
where $d=r+q\ge r$. Since $d'<0$, both $d^{-1}(1+b)$ and
$d^{-1}(1-b)=k$ are increasing. Their product is $M$, proving the
required logarithmic slope bound $m_0=0$.

It remains to check the matrix inequality in this tail. Write $N$ for the
matrix subtracted from $Q$ in that inequality, so that
\begin{equation}\label{cert:N}
 N=\begin{pmatrix}
 A_*b^2-\eta_*b+\nu_*/4+l^2/\alpha&\nu_*/2-\eta_*b\\
 \nu_*/2-\eta_*b&\nu_*
 \end{pmatrix}
 +\mathsf D\binom{\mathsf E}{cb}
                  \begin{pmatrix}\mathsf E&cb\end{pmatrix}.
\end{equation}
Equations \eqref{cert:left-ratios} give $\mathsf E<.05072$.
Using $b<.028$, $K_0<.56$, and $J_0<.72$ in \eqref{cert:N} gives
\[
 N_{11}<.224,\qquad N_{22}<.802,\qquad .34<N_{12}<.386.
\]
Since $Q_{11}=.43$, $Q_{12}=.13$, and $Q_{22}=2.72$, the leading
minor of $Q-N-\operatorname{diag}(.13,.75)$ is positive, and its
determinant is at least
\[
 (.43-.224-.13)(2.72-.802-.75)-(.13-.386)^2
 =.023232>0.
\]
This proves the local matrix bound on $a\le-6$.

\subsection*{The positive tail}
For $r=a\ge10$ put $u=r^{-2}$ and
$\theta=\log(r\sqrt{2\pi})$. Gaussian tail integration gives
\[
 F(u)=\int_0^\infty e^{-v-uv^2/2}\,dv,\qquad
 \Phi(-r)=\frac{\varphi(r)}r F(u),\qquad
 k=\frac r2+\frac{C(u,\theta)}r,
\]
where
\begin{equation}\label{cert:C}
 C=F\theta+G,\qquad
 G=\frac{F-1}{2u}-F\log F.
\end{equation}
Let $\mathcal U=u\partial_u$, with $\theta$ held fixed, and let
$\mathcal D=r\,d/dr=\partial_\theta-2\mathcal U$. Define
\[
 B=(1-\mathcal D)C,\qquad J=(2-\mathcal D)B,\qquad
 R=(3-\mathcal D)J.
\]
Here $B,J,R$ are auxiliary scalar functions used only in this appendix.
Differentiating $k$ gives
\begin{equation}\label{cert:right-derivatives}
 b=.5-uB,\qquad l=J/r^3,\qquad l'=-R/r^4.
\end{equation}
We first establish uniform bounds for these three derivatives:
\begin{equation}\label{cert:right-errors}
 |B-(\theta-1.5)|<.04,\qquad
 |J-(2\theta-4)|<.13,\qquad
 |R-(6\theta-14)|<.75.
\end{equation}

Differentiation under the integral gives, for $0\le u\le.01$,
\[
 1-u\le F\le1,\qquad |F'|\le1,\qquad
 |F''|\le6,\qquad |F'''|\le90,\qquad |F^{(4)}|\le2520.
\]
Indeed the $j$th derivative is bounded in absolute value by
$(2j)!/2^j$. To control the apparent quotient in $G$, use
$(F(u)-1)/u=\int_0^1F'(tu)\,dt$. Since $F\ge.99$ and
$|1+\log F|\le1$, differentiating \eqref{cert:C} yields
\[
 |G''|\le15+6+1/.99<23,\qquad
 |G'''|\le315+90+18/.99+1/.99^2<435.
\]
The values $F(0)=1$, $F'(0)=-1$, $G(0)=-1/2$, and $G'(0)=5/2$
therefore give
\[
 C=\theta-.5+(2.5-\theta)u+R_0,\qquad
 g=11.5+3\theta,\qquad c_0=435+90\theta,
\]
with
\begin{equation}\label{cert:Euler-remainder}
 |\mathcal U^jR_0|\le2^jg u^2+c_0u^3\mathbf1_{\{j=3\}}
       \qquad(0\le j\le3).
\end{equation}
The same bounds apply to $\partial_\theta R_0$ with $g=3,c_0=90$;
all higher $\theta$ derivatives vanish. For example,
$\mathcal U^2R_0=uR_0'+u^2R_0''$ and
$\mathcal U^3R_0=uR_0'+3u^2R_0''+u^3R_0'''$, so these statements
follow directly from Taylor's theorem and the derivative bounds just given.

Expanding the three operators in \eqref{cert:right-derivatives}, the
linear-in-$u$ errors are respectively
\[
 (8.5-3\theta)u,\qquad (37-12\theta)u,\qquad
 (197-60\theta)u.
\]
The errors contributed by $R_0$ are bounded respectively by
\begin{equation}\label{cert:operator-errors}
 (5g+3)u^2,\qquad(30g+33)u^2,\qquad
 (210g+321+8c_0u)u^2.
\end{equation}
For example the second operator is
$2-3\mathcal D+\mathcal D^2$ and the third is
$6-11\mathcal D+6\mathcal D^2-\mathcal D^3$; substituting
$\mathcal D=\partial_\theta-2\mathcal U$ into
\eqref{cert:Euler-remainder} gives exactly the displayed constants.

Here is a rational way to check uniformity on the unbounded interval.
Write $x=\log(r/10)\ge0$, so $u=.01e^{-2x}$ and
$\theta=\theta_0+x$, where $3.221<\theta_0<3.222$.
The terms in \eqref{cert:operator-errors} decrease with $x$. For the
linear-in-$u$ terms, maximize a linear function times $e^{-2x}$ at
its endpoint or stationary point, using $e^{-t}\le1/(1+t)$ for $t\ge0$.
More explicitly, set
\[
 \begin{aligned}
 g_+&=11.5+3(3.222),& c_+&=435+90(3.222),\\
 a_B&=3(3.222)-8.5,& a_J&=12(3.222)-37,\\
 a_R&=197-60(3.222).
 \end{aligned}
\]
Upper bounds for the total errors are
\[
 \begin{array}{c|l|c}
 &\text{rational upper bound}&\text{a larger rational}\\\hline
 B&\displaystyle\frac{.015}{2-2a_B/3}
                  +\frac{5g_++3}{10^4}&.024\\[2mm]
 J&\displaystyle\frac{.06}{2-a_J/6}
                  +\frac{30g_++33}{10^4}&.102\\[2mm]
 R&\displaystyle\max\left\{\frac{197-60(3.221)}{100},
                    \frac{.3}{2+a_R/30}\right\}
       +\frac{210g_++321}{10^4}+\frac{8c_+}{10^6}&.624
 \end{array}
\]
These prove \eqref{cert:right-errors}. The bound on $\theta_0$ itself
can be checked with the logarithm series
\[
 \log y=2\sum_{j=0}^{N}\frac{z^{2j+1}}{2j+1}+\mathcal R_N,\qquad
 z=\frac{y-1}{y+1},\qquad
 |\mathcal R_N|\le
 \frac{2|z|^{2N+3}}{(2N+3)(1-z^2)},
\]
applied to $\theta_0=\tfrac12\log(200\pi)$. For a completely
specified enclosure, use Machin's identity
$\pi=16\arctan(1/5)-4\arctan(1/239)$ and the alternating arctangent
sums through degrees $47$ and $49$. In the logarithm series first divide
the argument by a power of two to place it in $[1,2)$, and take $N=35$
for that argument and for $\log2$. These rational operations give
$3.221523626198<\theta_0<3.221523626199$.

Combining \eqref{cert:right-errors} with the preceding Taylor bounds
on $C$ gives the following estimates needed in the lemmas:
\begin{equation}\label{cert:right-consequences}
 1.7<R/J\le3,\quad .482<b<.5,\quad 0<l<.0026,\quad
 .5<k/r<.528,\quad .493<p<.5.
\end{equation}
To see the ratio bounds directly, $J>2\theta-4.13>0$,
$R-1.7J>2.6\theta-8.171>0$, and $3J-R>.86$.
For the bounds involving $k$ and $p$, we have
$C\ge .99\theta-.5>0$ by \eqref{cert:C}, whereas its Taylor expansion gives
\[
 C-(\theta-.5)\le
 u\{2.5-\theta+.01(11.5+3\theta)\}<0.
\]
Thus $0<C<\theta-.5<\theta-.475$, and $0<uCJ<.075$. In particular,
\[
 p=.5+u\{-B+(.5+uC)J\}-l^2/\alpha,\qquad
 -.605<-B+.5J<-.395,
\]
which proves $.493<p<.5$. The inverse-power majorants used for these consequences take their maxima at $r=10$,
as follows by differentiating the linear or quadratic polynomials in
$\log(r/10)$ times the indicated inverse power of $r$. For example,
$u(\theta-.475)(2\theta-3.87)$ decreases because
$1/(\theta-.475)+2/(2\theta-3.87)<2$ at $\theta=3.221$ and hence
for every larger $\theta$. Similarly $B/r<.18$, and consequently
\[
 |h_0|\le .528(.18)+2(.22)(.0026)/\alpha<.14<.67\sqrt p.
\]
These bounds prove the claimed estimates on $p,h_0$ and the $T_1$ row
of the coefficient table; for example
\[
 .367<K_0<.4,\qquad \mathsf D<D_*/.493<.640,\qquad
 J_0<D_*c(.5)/.493<.48,\qquad \mathsf L<.0026/\alpha<.20.
\]
They also give $-.3<z<0$ and
\[
 .8<1+3(.482)-3(.528)\le h_*
       \le1+3(.5)-1.7(.5)<1.77.
\]

To check the remaining estimates involving $d$, let $V$ have density
$F(u)^{-1}e^{-v-uv^2/2}$ on $v>0$. The Gaussian excess above $r$
is $V/r$, so
\[
 rd=\mathbb EV,\qquad 1-qd=r^{-2}\operatorname{Var}V.
\]
Taylor's lower bound in the numerator and $F\le1$ give
$\mathbb EV\ge1-3u$ and $\mathbb EV^2\ge2(1-6u)$.
Monotone weighting of the exponential density gives
$\mathbb EV\le1$. Hence $d\ge.97/r$,
$\operatorname{Var}V\ge.88$, and
\[
 \chi\le\frac{J}{.97r^2}<.027<.097.
\]
Equation \eqref{cert:derived} now implies
\[
 \frac{d\log M}{d\log q}
 \ge2(.88)-\frac{.027}{1-.5^2}>.68.
\]
This verifies the logarithmic slope required by Lemma~\ref{tr:scalar-basic}.

For the matrix inequality, use $\mathsf E\le p+.482(.75-.482)$
in \eqref{cert:N}. The first part of $N_{11}$ increases with $b$
on $[.482,.5]$, and $(p+v)^2/p$ increases with $p$ when
$0<v<p$. Therefore
\[
 \begin{split}
 N_{11}&\le A_*/4-\eta_*/2+\nu_*/4+.0026^2/\alpha
       +D_*\frac{(.5+.482(.75-.482))^2}{.5}<.403,\\
 N_{22}&\le\nu_*+D_*\frac{(c/2)^2}{.493}<1.12,
 \qquad .080<N_{12}<.13.
 \end{split}
\]
For the off-diagonal estimate use
$N_{12}=\nu_*/2+b(cK_0-\eta_*)$ and $.367<K_0<.4$.
Thus $Q-N-\operatorname{diag}(.022,.75)$ has positive leading minor
and determinant at least
\[
 (.43-.403-.022)(2.72-1.12-.75)-.05^2=.00175>0.
\]
This finishes the positive tail.

\subsection*{Polynomial enclosures on the central interval}
We next provide all data and rules for the finite certificate on $[-6,10]$.
On each row of Table~\ref{cert:centers}, write $a=w+rx$, $-1\le x\le1$.
The corresponding closed intervals are consecutive and cover $[-6,10]$.
The tabulated constants are rational seeds; we do not assume that they are
Gaussian function values. The differential residuals and the two analytic
anchors $q(10)$ and $k(-6)$ will certify the entire piecewise polynomial
family, including all its initial coefficients.
Construct polynomials $P,\widehat k$ of degree $14$ from their tabulated constant
terms by the recurrences
\begin{equation}\label{cert:recurrence}
 D=P-a,\qquad
 P_{i+1}=\frac r{i+1}(PD)_i,\qquad
 \widehat k_{i+1}=\frac r{i+1}(1-\widehat k D)_i\quad(0\le i<14).
\end{equation}
The subscript denotes a coefficient, so the constant $1$ contributes only
when $i=0$. At each step only coefficients already computed occur on the
right. The letters $P,\widehat k,D$ denote polynomial approximations to $q,k,d$,
respectively, throughout this central calculation.

\begin{table}[htbp]
\centering
\caption{Rational input data for the central polynomial certificate.}
\label{cert:centers}
\begin{tabular}{r r l l}
\hline
$w$&$r$&$P_0$&$\widehat k_0$\\\hline
$-5.25$&$.75$&$.0000004128471303$&$.1842076703035094$\\
$-3.75$&$.75$&$.0003526268606772$&$.2507500253088710$\\
$-2.5$&$.5$&$.0176378254869167$&$.3531628936050266$\\
$-1.5$&$.5$&$.1387897504588508$&$.4981884857033017$\\
$-.5$&$.5$&$.5091604338370335$&$.7246172144765372$\\
$.25$&$.25$&$.9635539794164037$&$.9475979387933716$\\
$1.25$&$.75$&$1.7288166273310539$&$1.3000948982566192$\\
$3$&$1$&$3.2830986549304364$&$2.0126493036044195$\\
$5$&$1$&$5.1865039671258417$&$2.9046537878842420$\\
$7$&$1$&$7.1375456132265036$&$3.8366560415481663$\\
$9$&$1$&$9.1085231050028685$&$4.7898159350739631$\\\hline
\end{tabular}
\end{table}

For a polynomial $A$, let $[A]$ denote the sum of the absolute values of
its coefficients. Direct use of \eqref{cert:recurrence} gives
\begin{equation}\label{cert:residuals}
 r[(PD)_{\deg\ge14}]<
 \begin{cases}2\cdot10^{-9},&w=-3.75,\\7\cdot10^{-11},&\text{otherwise},\end{cases}
 \qquad r[(\widehat k D)_{\deg\ge14}]<2\cdot10^{-10}.
\end{equation}
For each adjacent pair, the values of both $P$ and $\widehat k$ at their common
endpoint differ by less than $2\cdot10^{-10}$. Also
\begin{equation}\label{cert:additional-poly}
 D>.09,\qquad [\widehat k]<\begin{cases}2,&w\le.25,\\5.6,&w>.25,
 \end{cases}
\end{equation}
and the first-row value $\widehat k(-1)$ differs from $.16237766$ by less than
$3\cdot10^{-9}$, while the last-row value $P(1)$ differs from
$10.098093234$ by less than $2\cdot10^{-10}$.
These are rational coefficient checks. For example the first residual is
bounded by
\[
 r\sum_{i=14}^{28}\left|\sum_{j=0}^{14}P_jD_{i-j}\right|,
\]
with missing coefficients taken as zero. The bound on $D$ follows from
the quartic lower-bound rule \eqref{cert:Bernstein} below, after charging
discarded coefficients to the remainder.

Here are analytic anchors for those coefficient checks. For any $r>0$,
Taylor's formula under the integral gives alternating upper and lower
bounds for $I_0(r)$ by
\begin{equation}\label{cert:Mills}
 S_N(r)=\sum_{i=0}^{N}\frac{(-1)^i(2i-1)!!}{r^{2i+1}},
 \qquad S_{2j+1}(r)\le I_0(r)\le S_{2j}(r),
\end{equation}
where $(-1)!!=1$. At $r=10$, use $S_{11},S_{12}$ and $q(10)=1/I_0(10)$;
at $r=6$, use $S_{13},S_{14}$ and
$|k(-6)-I_0(6)|\le2\varphi(6)/36$. Exact substitution gives
\[
 |q(10)-10.098093234|<4\cdot10^{-10},\qquad
 |k(-6)-.16237766|<4\cdot10^{-8}.
\]
The rough bound $\varphi(6)<7\cdot10^{-9}$ used here follows, for example,
from $\sqrt{2\pi}>12/5$ and
$e^{18}>\sum_{i=0}^{59}18^i/i!$: division gives
$(5/12)/\sum_{i=0}^{59}18^i/i!<7\cdot10^{-9}$.

These anchors and residuals imply the following
uniform, rather than just endpoint, errors:
\begin{equation}\label{cert:global-errors}
 |q-P|\le e_q=10^{-8},\qquad |k-\widehat k|\le e_k=8\cdot10^{-7}.
\end{equation}
Indeed the equation for $q-P$ has multiplier $r(q+P-a)>0$, so backward
integration from $a=10$ is stable. Its absolute error increases by at most
twice each bound in \eqref{cert:residuals} and by each joining error.
The resulting bound is
\[
 6\cdot10^{-10}+2(2\cdot10^{-9}+10\cdot7\cdot10^{-11})
       +10\cdot2\cdot10^{-10}=8\cdot10^{-9}<e_q.
\]
The equation for $k-\widehat k$ has multiplier $-rd<0$, so forward integration
from $a=-6$ is stable. The additional forcing is bounded by $r[\widehat k]e_q$.
The portions with $[\widehat k]<2$ and $[\widehat k]<5.6$ have total lengths $6.5$ and
$9.5$, respectively. Thus an upper bound for its accumulated error is
\[
 4.3\cdot10^{-8}+22\cdot10^{-10}\cdot2
       +20\cdot10^{-10}
       +10^{-8}(6.5\cdot2+9.5\cdot5.6)
 =7.114\cdot10^{-7}<e_k.
\]

Construct the following \emph{full} polynomials before truncating any of
them:
\begin{equation}\label{cert:aux-polys}
 B=1-\widehat k D,\qquad J=\widehat k-P+aB,\qquad Y=2B-PD+aJ.
\end{equation}
They approximate $b,l,l'$. Their error bounds, obtained directly from
\eqref{cert:identities} and \eqref{cert:global-errors}, are
\begin{equation}\label{cert:aux-errors}
 \begin{split}
 E_b&=[\widehat k]e_q+e_k(e_q+[D]),\\
 E_l&=e_q+e_k+[a]E_b,\\
 E_{l'}&=2E_b+([P]+[D]+e_q)e_q+[a]E_l.
 \end{split}
\end{equation}
The order in \eqref{cert:aux-polys} matters for the strength of the
certificate: form $B,J,Y$ with the full degree-$14$ inputs, and only then
apply the quartic arithmetic described next.

\subsection*{The rational positivity certificate}
An enclosure $(A,e)$ means that the function in question differs from
$A(x)$ by at most $e$ on $[-1,1]$. Replace a polynomial by its terms of
degree at most four and add the sum of the absolute values of all discarded
coefficients to $e$. Addition and subtraction add error bounds. For a
product use the coefficient convolution and the error
\begin{equation}\label{cert:product-error}
 [A]e_B+[B]e_A+e_Ae_B,
\end{equation}
then truncate and charge the discarded coefficients as before. Scalar
multiplication scales the error by the absolute value of the scalar.
Initialize this arithmetic with
\[
 (q,k,d,b,l,l')\longleftrightarrow
 (P,\widehat k,D,B,J,Y),\qquad
 (e_q,e_k,e_q,E_b,E_l,E_{l'}),
\]
using the full polynomials in \eqref{cert:aux-polys}. Truncate each of
these six initial enclosures to degree four, charging discarded coefficients
to its error, before evaluating the groups below.
For a resulting quartic $A(x)=\sum_{j=0}^4A_jx^j$, a rational lower bound is
\begin{equation}\label{cert:Bernstein}
 \min_{\substack{\sigma=\pm1\\0\le i\le4}}
       \sum_{j=0}^i\frac{\binom ij}{\binom4j}\sigma^jA_j-e.
\end{equation}
This is the Bernstein coefficient bound on each of $[0,1]$ and $[-1,0]$:
on $x=\sigma t$ the displayed coefficients are precisely the coefficients
in the degree-four Bernstein basis, whose basis functions are nonnegative
and sum to one.

We list every polynomial to which this rule is applied. In a given row,
use the corresponding bin in Lemma~\ref{sc:bounds} to define the bounds
\[
 b_{\min},\ b_{\max},\ \mathsf L_{\max},\ \mathsf D_{\max},\qquad
 K_{\min},\ K_{\max},\ J_{\max},\ \ell_1.
\]
The constant $m_0$ is from Lemma~\ref{tr:scalar-basic}.
The triples $(z_{\min},z_{\max},\chi_{\max})$ are
$(-.08,.70,.070)$ for rows to the left of $-1$, and
$(-.475,0,.097)$ for rows to the right. Within the quartic arithmetic set
\[
 \begin{aligned}
 p&=b+kl-l^2/\alpha,&
 \mathsf E&=p+b(.75-b)-l^2/\alpha,\\
 H&=l+3bl+kl',& f_1&=cb,
 \end{aligned}
\]
\[
 \begin{aligned}
 h&=.43-\ell_1-A_*b^2+\eta_*b-\nu_*/4-l^2/\alpha,\\
 v&=1.97-\nu_*,\qquad o=.13-\nu_*/2+\eta_*b.
 \end{aligned}
\]
The four groups are as follows:
\begin{align*}
 \mathcal G_1:\quad& b-b_{\min},\ b_{\max}-b,\ l,\
       \alpha\mathsf L_{\max}-l,\\
 &p(7+a^2)-1\quad\text{on rows with }w<0,\\
 &p-.25\quad\text{on rows with }0<w<4,\qquad
   p-.46\quad\text{on rows with }w>4,\\
 &.67^2p-\{k(b-.5)-2(b-.28)l/\alpha\}^2;\\[1mm]
 \mathcal G_2:\quad& l'-z_{\min}l,\ z_{\max}l-l',\
       H-.8l,\ 1.77l-H,\ \chi_{\max}d-l,\\
 &2(1-qd)(1-b^2)-2bld-m_0d^2(1-b^2);\\[1mm]
 \mathcal G_3:\quad& p-D_*/\mathsf D_{\max}\quad
                         \text{(omitted for bin }O\text{)},\\
 &D_*\mathsf E-K_{\min}p,\ K_{\max}p-D_*\mathsf E,\
       J_{\max}p-D_*f_1;\\[1mm]
 \mathcal G_4:\quad& hp-D_*\mathsf E^2,\\
 &p(hv-o^2)-D_*(\mathsf E^2v-2\mathsf E\,o\,f_1+h f_1 f_1).
\end{align*}
Products are evaluated from left to right, with repeated multiplication
for powers. In the last expression the product in the middle is
$2\mathsf E\,o\,f_1$. The bounds below are the minimum over all members
of a group, so positivity proves every inequality listed in it.

\begin{table}[htbp]
\centering
\caption{Lower bounds for the four polynomial groups, in units of $10^{-3}$.}
\label{cert:group-bounds}
\begin{tabular}{r r r r r}
\hline
$w$&$\mathcal G_1$&$\mathcal G_2$&$\mathcal G_3$&$\mathcal G_4$\\\hline
$-5.25$&$7$&$1$&$1$&$1$\\
$-3.75$&$1$&$1$&$2$&$2$\\
$-2.5$&$2$&$2$&$2$&$5$\\
$-1.5$&$2$&$2$&$2$&$8$\\
$-.5$&$2$&$3$&$5$&$14$\\
$.25$&$2$&$2$&$8$&$11$\\
$1.25$&$2$&$.8$&$8$&$20$\\
$3$&$2$&$.8$&$8$&$7$\\
$5$&$1$&$.8$&$10$&$15$\\
$7$&$3$&$.6$&$11$&$11$\\
$9$&$2$&$.35$&$12$&$9$\\\hline
\end{tabular}
\end{table}

Table~\ref{cert:group-bounds} is obtained by the rational operations
\eqref{cert:recurrence}, \eqref{cert:aux-polys},
\eqref{cert:aux-errors}, \eqref{cert:product-error}, and
\eqref{cert:Bernstein}, with no additional function evaluations.
The accompanying checker reproduces all $44$ bounds, all $22$
differential residual bounds, the $20$ joining bounds, the anchor checks,
and the coefficient norm bounds used to obtain
\eqref{cert:global-errors}. The same endpoint enclosures give
$q(-1)>.287$, $q(-2)<.056$, $q(-3)<.005$, and $q(4)>4$, as used
in the secant argument. The logarithm series above also certifies
\[
 \log\frac{q(-1)}{q(-2)}>1.62,\quad
 \log\frac{q(-1)}{q(-3)}>4,\quad
 \log\frac{q(4)}{q(-2)}>4,\quad
 \log\frac{q(4)}{q(-3)}>6.5.
\]

We explain why these finite checks finish the proof. Group $\mathcal G_1$
gives $0<b<1/2$, $l>0$, the bounds for $p,h_0$, and the bounds for
$b,\mathsf L$. Its strict positive margins give $p>0$.
Group $\mathcal G_2$ gives the bounds for $z,h_*,\chi$ after division by
$l$ or $d$, and the logarithmic slope bound by division by
$d^2(1-b^2)$. Group $\mathcal G_3$ gives the remaining coefficient bounds
after division by $p$; the $O$ bound for $\mathsf D$ instead follows from
$D_*<.295$ and the lower bound on $p$.
Finally
\[
 Q-N-\operatorname{diag}(\ell_1,.75)
 =\begin{pmatrix}h&o\\o&v\end{pmatrix}
     -\frac{D_*}{p}\binom{\mathsf E}{f_1}
                             \begin{pmatrix}\mathsf E&f_1\end{pmatrix}.
\]
Its leading minor and determinant, multiplied by $p$, are exactly the
two expressions in $\mathcal G_4$. Their positivity proves the matrix
inequality. Continuity supplies every bin endpoint, including the stronger
of the two adjacent logarithmic slope bounds. Together with the two tails,
this proves both lemmas on the whole real line.
\end{proof}
